% from ega3-7-fr.tex
\oldpage[III]{198}
Nous allons considérer le $\partial$-foncteur
$\mathscr M\leadsto\mathscr T_\bullet(\mathscr M)$ (aussi noté
$\mathscr T_\bullet(\mathscr P_\bullet,\mathscr M)$) dans la catégorie des
$\mathscr O_Y$-Modules quasi-cohérents, à valeurs dans la catégorie des
$\mathscr O_Y$-Modules quasi-cohérents (en vertu de (6.2.3)), défini par
\[
  \mathscr T_n(\mathscr P_\bullet,\mathscr M)
  =\mathscr T_n(\mathscr M)
  =\mathscr H^{-n}
    \bigl(f,\mathscr P^\bullet\otimes_{\mathscr O_Y}\mathscr M\bigr)
  \qquad\text{pour }n\in\mathbf Z,
  \tag{7.7.1.1}\label{III.7.7.1.1.fr}
\]
où $\mathscr P^\bullet$ est le complexe dont le terme de degré $j$ est
$P_{-j}$, l'opérateur de dérivation étant donc alors de degré $+1$. Le foncteur
$\mathscr T_\bullet$ ainsi défini est un \emph{foncteur homologique} en
$\mathscr M$ (6.2.6).
\end{env}

\begin{env}[7.7.2]
Soit $g:Y'\to Y$ un morphisme, et posons
$X'=X_{(Y')}=X\times_Y Y'$ et $f'=f_{(Y')}:X'\longrightarrow Y'$,
qui est un morphisme quasi-compact et séparé; soit d'autre part
$\mathscr P'_\bullet=\mathscr P_\bullet\otimes_{\mathscr O_Y}\mathscr O_{Y'}$;
c'est un complexe de $\mathscr O_{X'}$-Modules quasi-cohérents qui sont
$Y'$-\emph{plats} en vertu de (\textbf{I}, 9.1.12) et ($\mathbf{0}_{\mathrm I}$, 6.2.1). Nous poserons
(avec les mêmes conventions sur les degrés)
\[
  \mathscr T^{Y'}_\bullet(\mathscr M')
  =\mathscr H^\bullet
    \bigl(f',\mathscr P'^\bullet\otimes_{\mathscr O_{Y'}}\mathscr M'\bigr)
  =\mathscr H^\bullet
    \bigl(f',\mathscr P^\bullet\otimes_{\mathscr O_Y}\mathscr M'\bigr)
  \tag{7.7.2.1}\label{III.7.7.2.1.fr}
\]
qui est un foncteur homologique en le $\mathscr O_{Y'}$-Module quasi-cohérent
$\mathscr M'$. Lorsque $Y'$ est un schéma affine d'anneau $A'$, on écrira
$\mathscr T^{A'}_\bullet$ au lieu de $\mathscr T^{Y'}_\bullet$; pour tout
$A'$-module $M'$, on a alors
$\mathscr T^{A'}_\bullet(\widetilde{M}')=\bigl(\Gamma(Y',\mathscr T^{Y'}_\bullet(\widetilde{M}'))\bigr)^\sim$;
on posera
$T^{A'}_\bullet(M')=\Gamma(Y',\mathscr T^{Y'}_\bullet(\widetilde{M}'))$,
qui est un foncteur homologique de $A'$-modules, à valeurs dans la catégorie des
$A'$-modules.
On observera que si $Y=\operatorname{Spec}(A)$ est aussi affine, le foncteur de
$A'$-modules $T^{A'}_\bullet$ coïncide avec le foncteur obtenu par extension des
scalaires de $A$ à $A'$ à partir du foncteur homologique de $A$-modules
$T^A_\bullet$ (7.1.3) : en effet, soit $g:Y'\to Y$ le morphisme correspondant à
l'homomorphisme d'anneaux $A\to A'$, et soit $g':X'\to X$ le morphisme
correspondant qui est affine (\textbf{II}, 1.6.2); si $\mathfrak U$ est un recouvrement
ouvert affine de $X$, $\mathfrak U'=g'^{-1}(\mathfrak U)$ est un recouvrement
ouvert affine de $X'$; en vertu de (6.2.2), tout revient à voir que
$C^\bullet\bigl(\mathfrak U,\mathscr P_\bullet\otimes_{\mathscr O_Y}g_*(\mathscr M')\bigr)=C^\bullet\bigl(\mathfrak U',\mathscr P_\bullet\otimes_{\mathscr O_Y}\mathscr M'\bigr)$,
et finalement, que pour tout ouvert affine $U$ de $X$, en posant
$U'=g'^{-1}(U)$, on a
$\Gamma\bigl(U,\mathscr P_\bullet\otimes_{\mathscr O_Y}g_*(\mathscr M')\bigr)=\Gamma\bigl(U',\mathscr P_\bullet\otimes_{\mathscr O_Y}\mathscr M'\bigr)$,
ce qui est trivial (\textbf{I}, 1.3 et 3.2).

En particulier, si $U$ est un ouvert de $Y$, on a, pour tout
$\mathscr O_Y$-Module quasi-cohérent $\mathscr M$
\[
  \mathscr T^U_\bullet(\mathscr M|U)
  =\bigl(\mathscr T_\bullet(\mathscr M)\bigr)|U.
  \tag{7.7.2.2}\label{III.7.7.2.2.fr}
\]
\end{env}

\begin{env}[7.7.3]
Pour tout $\mathscr O_Y$-Module quasi-cohérent $\mathscr M$, on a un
homomorphisme canonique, fonctoriel en $\mathscr M$ :
\[
  \mathscr T_p(\mathscr O_Y)\otimes_{\mathscr O_Y}\mathscr M
  \longrightarrow
  \mathscr T_p(\mathscr M).
  \tag{7.7.3.1}\label{III.7.7.3.1.fr}
\]

En effet, si $Y$ est affine, cet homomorphisme a été défini en (7.2.2); cette
définition s'étend sans peine au cas général, en remarquant que si $U$, $V$ sont
deux ouverts affines de $Y$ tels que $V\subset U$, le diagramme
